\documentclass[11pt]{amsart}
\usepackage[T1]{fontenc}
\usepackage[utf8]{inputenc}
\usepackage{amssymb}
\usepackage[margin=1.05in]{geometry}
\usepackage[hidelinks]{hyperref}
\newcommand{\stamptime}{\the\numexpr\time/60\relax:\ifnum\numexpr\time-60*(\time/60)\relax<10 0\fi\the\numexpr\time-60*(\time/60)\relax}
\AddToHook{shipout/foreground}{\put(0.5\paperwidth,-\paperheight+0.55cm){\makebox[0pt]{\tiny\ttfamily unit/doc/\jobname.tex \ --- \ compiled \today, \stamptime}}}
\newtheorem{theorem}{Theorem}
\newtheorem{proposition}[theorem]{Proposition}
\newtheorem{lemma}[theorem]{Lemma}
\newtheorem{corollary}[theorem]{Corollary}
\theoremstyle{definition}
\newtheorem{definition}[theorem]{Definition}
\newtheorem{statement}[theorem]{Statement}
\theoremstyle{remark}
\newtheorem{remark}[theorem]{Remark}
\newcommand{\Om}{\Omega}
\newcommand{\code}[1]{\texttt{#1}}
\newcommand{\cT}{\mathsf{isTrunc}}
\newcommand{\cC}{\mathsf{isConn}}
\newcommand{\refl}{\mathsf{refl}}
\newcommand{\ap}{\mathsf{ap}}
\newcommand{\Fib}{\mathrm{Fib}}
\newcommand{\fib}{\mathrm{fib}}
\newcommand{\tr}[1]{\|#1\|}
\newcommand{\Ab}{\mathrm{Ab}}
\newcommand{\Tor}{\mathsf{Tor}}
\newcommand{\SCt}{\mathsf{SC3}}
\newcommand{\SCf}{\mathsf{SC4}}
\newcommand{\Tri}{\mathsf{Triples}}
\newcommand{\Trit}{\mathsf{Triples3}}
\newcommand{\Del}{\mathsf{Deloopable}}
\newcommand{\Unloop}{\mathsf{Unloop}}
\newcommand{\real}{\mathsf{realise}}
\newcommand{\pteq}{\simeq_*}
\newcommand{\Map}{\mathrm{Map}_*}
\newcommand{\PtEq}{\mathrm{PtEquiv}}
\newcommand{\lp}{\mathsf{loop}}
\newcommand{\Z}{\mathbb{Z}}
\newcommand{\U}{\mathcal{U}}
\title{Delooping as an embedding}
\date{28 October 2026}
\begin{document}
\maketitle

\noindent\textbf{Purpose.} Informal proof for \code{doc/delooping-embedding-handoff.txt}:
looping is an embedding of pointed $2$-connected $4$-types into pointed
$1$-connected $3$-types, the resulting description of the former inside the
classification by triples, the delooping of $R=\Fib\bar k$, and the precise
statement of the marked universal property (step 2 of the handoff) together with
an analysis of what the embedding does and does not give for it.

\noindent\textbf{Attribution.} The organising principle (delooping is an
embedding, and $\pi_4(S^3)$ is read off inside the deloopable part of the
classification) and the sphere universal-property strategy are David Wärn's;
the delooping machinery (\code{DeloopMaps}, DW1 \S3) is his. The classification
of pointed simply connected $3$-types by triples is Hart--Milner's (Owen
Milner's equivalence; \code{doc/postnikov-classification.pdf}).

\noindent\textbf{Conventions.} As in the development, $\cC\,n\,X$ means that $X$
is $(n-2)$-connected and $\cT\,n\,X$ that $X$ is an $(n-2)$-type. Write
\[
  \SCt = (X:\U)\times(x:X)\times\cC\,3\,X\times\cT\,5\,X,\qquad
  \SCf = (Y:\U)\times(y:Y)\times\cC\,4\,Y\times\cT\,6\,Y,
\]
(pointed $1$-connected $3$-types, \code{ClassPost.SC3}; pointed $2$-connected
$4$-types). Both are of the form $\mathsf{PP}_P=(X:\U)\times(x:X)\times P\,X$
with $P$ a family of propositions, so \code{ClassBase.PtProp} applies:
\begin{equation}\label{eq:ptprop}
  (u=v)\ \simeq\ \PtEq(u,v),\qquad
  (v:\mathsf{PP}_P)\times\PtEq(u,v)\ \text{is contractible},
\end{equation}
the first map being transport of $\mathrm{id}_u$ (the fundamental theorem of
identity types applied to the second). Here
$\PtEq(u,v)=(F:\Map(u,v))\times\mathsf{isEquiv}\,F$. All of $\SCt$, $\SCf$, $\Tri$ and the types below over them live in
$\U_{\ell+1}$ (the carriers in $\U_\ell$), so the fundamental theorem,
\code{sigma-base} and \code{emb-fib-prop} are used at level $\ell+1$; $\Tor_R$
below is in $\U_\ell$. Looping is
\[
  \Om:\SCf\to\SCt,\qquad \Om(Y,y,c,t)=(y=y,\ \refl_y,\ c',\ t'),
\]
with $c'=\code{conn-id}$ and $t'=t\,y\,y$. For pointed maps
$\lp:\Map(Y,Y')\to\Map(\Om Y,\Om Y')$ is \code{DeloopMaps.Maps.loop}; its fibre
over $f$ is (up to the symmetry of $=$) $\Unloop f=(F:\Map(Y,Y'))\times(f=\lp F)$.

\section{Looping is an embedding}

\begin{lemma}[looping of maps]\label{lem:loopmaps}
Let $Y$ be pointed with $\cC\,4\,Y$ and $Y'$ pointed with $\cT\,6\,Y'$. Then
$\lp:\Map(Y,Y')\to\Map(\Om Y,\Om Y')$ is an equivalence.
\end{lemma}
\begin{proof}
DW1 Corollary 11 (\code{Over.cor11}) at $c=2$: the hypotheses are
$\cC\,(c+2)\,Y$ and $\cT\,((c+1)+(c+1))\,Y'$, i.e.\ $\cC\,4$ and $\cT\,6$;
then each $\Unloop f$ is contractible, and a map with contractible fibres is
an equivalence. This is \code{SourceS3.LoopEquiv} at $c=2$ (Lemma loopeq3 of
\code{doc/br-universal.pdf}); its inverse sends $f$ to the centre of
$\Unloop f$, and $\lp(\text{centre})=f$ is the second component.
\end{proof}

\begin{lemma}[looping of equivalences]\label{lem:loopeq}
For $u=(Y,y,..)$, $v=(Y',y',..)$ in $\SCf$, write $\Om_*:\PtEq(u,v)\to
\PtEq(\Om u,\Om v)$ for $(F,e)\mapsto(\lp F,\ \code{loop-pe-equiv}\,F)$
(\code{LoopIso.loopPE}). It has a section $\delta$:
$\Om_*(\delta g)=g$ for every $g$.
\end{lemma}
\begin{proof}
For $g=(f,e)$ let $F$ be the inverse of Lemma~\ref{lem:loopmaps} at $f$, so
$\lp F=f$. Since $\lp F=f$ is an equivalence and $Y$, $Y'$ are $0$-connected
($\cC\,2$, from $\cC\,4$ by \code{conn-down}), $F$ is an equivalence
(Whitehead's principle, \code{Whitehead.Pointed.equiv}). Put
$\delta g=(F,\ldots)$. Then $\Om_*(\delta g)=(\lp F,\ldots)=(f,e)$: the first
components agree by Lemma~\ref{lem:loopmaps}, the second ones lie in the
proposition $\mathsf{isEquiv}\,f$.
\end{proof}

\noindent Only the section is used: a retract of a contractible type is
contractible, so we never need that $\Om_*$ is an equivalence, nor that it is
compatible with $\ap_\Om$ under~(\ref{eq:ptprop}).

\begin{theorem}[E]\label{thm:E}
For every $u:\SCf$, the type $(v:\SCf)\times(\Om u=\Om v)$ is contractible.
Consequently $\ap_\Om:(u=v)\to(\Om u=\Om v)$ is an equivalence for all $u,v$
($\Om$ is an embedding), and every fibre $\fib_\Om(w)=(v:\SCf)\times(\Om v=w)$,
$w:\SCt$, is a proposition.
\end{theorem}
\begin{proof}
Fix $u$. Consider the three total spaces over $v:\SCf$
\[
  T_1=(v:\SCf)\times\PtEq(u,v),\quad
  T_2=(v:\SCf)\times\PtEq(\Om u,\Om v),\quad
  T_3=(v:\SCf)\times(\Om u=\Om v).
\]
(1) $T_1$ is contractible by~(\ref{eq:ptprop}) for $\SCf$
(\code{PtProp.contr}).

(2) $T_2$ is a retract of $T_1$: $s(v,g)=(v,\delta g)$, $r(v,G)=(v,\Om_*G)$,
$r(s(v,g))=(v,g)$ by Lemma~\ref{lem:loopeq}. Hence $T_2$ is contractible.

(3) $T_3$ is a retract of $T_2$: by~(\ref{eq:ptprop}) for $\SCt$ at $\Om u$ the
map $\mathsf{tr}_v:(\Om u=\Om v)\to\PtEq(\Om u,\Om v)$ is an equivalence with
inverse $\mathsf{un}_v$; take $s(v,p)=(v,\mathsf{tr}_v p)$, $r(v,g)=(v,\mathsf{un}_v g)$.
Hence $T_3$ is contractible.

By the fundamental theorem of identity types (\code{FundThm.FT} at the family
$v\mapsto(\Om u=\Om v)$ and the point $\refl_{\Om u}$), transport of
$\refl_{\Om u}$, $p\mapsto p_*\refl_{\Om u}$, is an equivalence
$(u=v)\to(\Om u=\Om v)$; it is $\ap_\Om$ (by $J$ on $p$: both are $\refl$ at
$\refl$). An embedding has propositional fibres (\code{Groupoid.emb-fib-prop}).
\end{proof}

\begin{remark}
The proof uses about $\Om$ only that it acts on $\PtEq$ through $\lp$ and that
$\lp$ has a section on equivalences. The hypotheses $\cC\,4$, $\cT\,6$ enter
only in Lemma~\ref{lem:loopmaps} (and $\cC\,2$ in Whitehead's principle).
\end{remark}

\section{The deloopable triples}

Recall $\Tri=(A\ B:\Ab)\times\Map(K(A,2),K(B,4))$ and the equivalence
$\real:\Tri\simeq\SCt$, $\real(A,B,k)=(\Fib k,(\mathrm{pt},k_0),..)$
(\code{PostnikovClassification.classification}).

\begin{definition}
$\Del\,t=\fib_\Om(\real\,t)=(v:\SCf)\times(\Om v=\real\,t)$.
\end{definition}

\begin{corollary}[C]\label{cor:C}
$\Del\,t$ is a proposition, and $(t,(v,p))\mapsto v$ is an equivalence
\[
  (t:\Tri)\times\Del\,t\ \simeq\ \SCf .
\]
\end{corollary}
\begin{proof}
A proposition by Theorem~\ref{thm:E}. The map is the composite of
$(t,z)\mapsto(\real\,t,z)$, an equivalence $(t:\Tri)\times\fib_\Om(\real t)\simeq
(w:\SCt)\times\fib_\Om(w)$ (Sigma over an equivalence of bases,
\code{SphereAxiom.sigma-base}, as in \code{PostnikovClassification.base}), and
$(w,(v,p))\mapsto v$, $(w:\SCt)\times\fib_\Om(w)\simeq\SCf$ (the total space of the
fibres of a map is its domain; inverse $v\mapsto(\Om v,(v,\refl))$, the
composite on the other side by $J$ on $p$).
\end{proof}

\noindent So pointed $2$-connected $4$-types are exactly the triples $(A,B,k)$
together with the \emph{property} that $\Fib k$ deloops. In the classification a
$4$-type carries no data beyond its loop space.

\section{Step 1: $R$ deloops}

Let $\bar t=(\Z,\Z/2,\bar k)$ and $R=\Fib\bar k=\real(\bar t)$
(\code{DeloopR.TR}).

\begin{proposition}
$\Del\,\bar t$, witnessed by $BR=\Fib\hat F$ (\code{DeloopR}).
\end{proposition}
\begin{proof}
$(BR,\mathrm{br}_0)$ lies in $\SCf$ by \code{BR-conn}, \code{BR-trunc};
\code{omega-BR} is a pointed equivalence $\Om BR\pteq R$, which
by~(\ref{eq:ptprop}) for $\SCt$ is a path $\Om(BR)=\real\,\bar t$.
\end{proof}

\begin{corollary}[uniqueness of the delooping of $R$]
$BR\pteq\Tor_R$, where $\Tor_R$ is the torsor delooping of
\code{DeloopR} (Corollary 19 of \code{doc/cross-term.pdf}).
\end{corollary}
\begin{proof}
$\Tor_R$ is $\cC\,4$ and $\cT\,6$ (\code{tor-conn}, \code{tor-trunc}) and
$\Om\Tor_R\pteq R$ (\code{tor-loops}); so it also inhabits $\Del\,\bar t$. Two
elements of a proposition are equal; the first projection gives
$BR=\Tor_R$ in $\SCf$, and~(\ref{eq:ptprop}) for $\SCf$ a pointed equivalence.
\end{proof}

\section{Step 2: the precise statement}

Under the $S^3$ axiom (segment parameters $S^3,\mathrm{base},\mathrm{surf},
\mathrm{ump}_3$), \code{SourceS3.Final} defines a class
$s\in[S^3,BR]_*$ (no representative is chosen). For $Y$ pointed write
$s^*:[BR,Y]_*\to[S^3,Y]_*$, $c\mapsto s;c$ (\code{HomotopyClasses.compC}).

\begin{statement}[M, the marked universal property of $BR$]
For every $(Y,y,..):\SCf$, $s^*:[BR,Y]_*\to[S^3,Y]_*$ is a bijection.
\end{statement}

\noindent For $T=\tr{S^3}_4$ (a pointed $2$-connected $4$-type,
\code{Sphere3Trunc}) and $BR$ this is exactly what step 3 consumes: the
two-map Yoneda argument \code{Pi4Sphere.Two4} needs $s^*$ bijective into
$P=BR$ and into $T$, and gives $\pi_4S^3\cong\pi_4T\cong\pi_4BR\cong\Z/2$. Into
$BR$ it is (a) below directly ($BR=\Fib\hat F$ literally); so M is consumed only
at $T$, and M for all of $\SCf$ is stronger than what is used.

\noindent\textbf{What is proved.}
\begin{enumerate}
\item[(a)] For $Y=\Fib K$, $K:K(A,3)\to_*K(B,5)$ any level-$3$ triple:
  \code{SourceS3.Final.rep-fib}; and along pointed equivalences,
  \code{rep-pteq}. No presentation is used.
\item[(b)] For every $Y:\SCf$: \code{Presentation4.Pres4} gives a level-$3$
  triple $K$ and $\Fib K\simeq Y$, pointed merely ($Y$ is connected); M is a
  proposition, so (a) applies. This is \code{Final4.rep-T} with $T$ replaced by
  an arbitrary $Y$.
\end{enumerate}
\begin{remark}[looped form; not used, not formalised] Let $e=\code{omega-BR}:\Om BR\pteq R$ and
\[
  r=(\Om s);[e]\ \in\ [\Om S^3,R]_*,
\]
a class ($\Om$ acts on classes; $s$ has no chosen representative). Since
$S^3$ and $BR$ are $\cC\,4$ (\code{Sphere3Trunc.Simply.s3-conn}, \code{BR-conn})
and $Y$ is $\cT\,6$, Lemma~\ref{lem:loopmaps} gives
bijections $[BR,Y]_*\cong[\Om BR,\Om Y]_*\cong[R,\Om Y]_*$ (the second by
$[e]^{-1}$) and $[S^3,Y]_*\cong[\Om S^3,\Om Y]_*$, and looping commutes with
composition and descends to classes (\code{Sphere3Trunc.loopC}, \code{loopC-comp}); so M at $Y$ is equivalent to the bijectivity of
\[
  r^*:[R,\Om Y]_*\longrightarrow[\Om S^3,\Om Y]_*,\qquad c\mapsto r;c .
\]
By Corollary~\ref{cor:C} the targets $\Om Y$ are exactly the $\real\,t$ with
$t$ deloopable. The existing results give this at once for the $t$ which are
loopings of level-$3$ triples (looping (a)).
\end{remark}

\section{The crux: are all deloopable triples loopings?}

For a level-$3$ triple $(A,B,K)$ let
\[
  \Om_3(A,B,K)=(A,B,\ e_2^{-1};\lp K;e_4)\ :\ \Tri,
\]
with $e_2:\Om K(A,3)\pteq K(A,2)$, $e_4:\Om K(B,5)\pteq K(B,4)$ the
Eilenberg--MacLane identifications, oriented as in the development
(\code{H4Bilinear.Lev.eA}, \code{eB}; \code{DeloopR}: $f=e_2;\bar k;e_4^{-1}$). For
$t=(A,B,k)$ put $f_t=e_2;k;e_4^{-1}:\Om K(A,3)\to_*\Om K(B,5)$.

\begin{proposition}\label{prop:crux}
For $t=(A,B,k):\Tri$:
\begin{enumerate}
\item $\Unloop f_t$ is a proposition, and it holds iff for every
  $p:\Om K(A,3)$, $f_t=D_p$ (the translate of $|f_t|$ by $p$,
  \code{DeloopMaps.Maps.D}); no coherence at $p=\refl$ is needed.
\item $\Unloop f_t\to\Del\,t$: if $f_t=\lp K$ then $\Fib K$ deloops $\real\,t$.
\item $\Del\,t\to\Unloop f_t$ holds if every pointed $2$-connected $4$-type is
  merely of the form $\Fib K$ (Pres4).
\end{enumerate}
Unconditionally, (1) and (2) with Theorem~\ref{thm:E} give that
$K\mapsto\Fib K$ is an embedding $\Trit\hookrightarrow\SCf$ (below); the Crux is
exactly its surjectivity. With (3), $\Del\,t\simeq\Unloop f_t$ and
\[
  \SCf\ \simeq\ (t:\Tri)\times\Unloop f_t\ \simeq\ \Trit
  =(A\ B:\Ab)\times\Map(K(A,3),K(B,5)),
\]
the classification one level up, with the inverse $K\mapsto\Fib K$.
\end{proposition}
\begin{proof}
(1) DW1 Corollary 10 at $c=2$, $t=1$: $K(A,3)$ is $\cC\,4$ and $K(B,5)$ is
$\cT\,7=\cT((2+1)+(2+1)+1)$, so each cell $(b:K(B,5))\times C\,a\,b$ and
$\Unloop f_t$ itself are propositions (\code{Over.Levels[2,1]}). If
$f_t=D_p$ for all $p$, then $(b_0,|f_t|,p\mapsto\ldots)$ inhabits the cell at the
base point, every cell is inhabited by connectedness (\code{conn-elim}), and
\code{lemma9.r} gives $\Unloop f_t$: this is \code{DeloopR.unloop-contr} (cells
\code{CrossTerm.Cross.cross-cell}; ``no pointing in $p$ is used''). Conversely
$f_t=\lp F$ gives the cells by \code{lemma9}, and $(b,h,c)$ in the cell at the
base point gives $f_t=D_p$ after the singleton contractions of \code{Levels}.
(\code{lemma9} is only a retract; the equivalence holds because both sides
are propositions.)

(2) $\Om\Fib K\pteq\Fib(\lp K)$ (\code{Fibres.omega-fib}) $\pteq\Fib f_t$ (along
$f_t=\lp K$) $\pteq\Fib k$ (along $e_2$, $e_4$, from the square
$f_t;e_4^{-1}=e_2^{-1};k$): the proof of \code{DeloopR.omega-BR} with $(Z,\Z/2,\hat
F)$ replaced by $(A,B,K)$. $\Fib K$ is $\cC\,4$ and $\cT\,6$ as in
\code{BR-trunc} (\code{MapBR.Triple3.trunc6}) and \code{BR-conn}: for a
general $K$ the connectivity is \code{conn-from-point} with
$\Om\Fib K\pteq\Fib k$ and \code{Triple.conn3}, as in \code{BR-conn}
(\code{Triple3} has no such field).

(3) $\Unloop f_t$ is a proposition, so we may take $Y$ deloops $\real\,t$ and
$Y\pteq\Fib K$. Then $\real\,t=\Om Y=\Om\Fib K=\real(\Om_3(A',B',K))$ by (2), so
$t=\Om_3(A',B',K)$ ($\real$ is an equivalence); transporting, $f_t$ becomes
$e_2;(e_2^{-1};\lp K;e_4);e_4^{-1}=\lp K$.

The embedding: $\Om(\Fib K)=\real(\Om_3 K)$ by (2), $\real\circ\Om_3$ has
propositional fibres ($\real$ is an equivalence and the fibre of $\Om_3$ over
$t$ is $\Unloop f_t$ up to the contractible choice of $k$, a proposition by
(1)), and $\Om$ is an embedding; if $g\circ h$ and $g$ are embeddings so is $h$.

For the last equivalence: $(t:\Tri)\times\Unloop f_t$ is
$(A\,B)\times(k)\times(K)\times(f_{(A,B,k)}=\lp K)$; for fixed $K$,
$k\mapsto f_{(A,B,k)}$ is an equivalence (pre- and post-composition with pointed
equivalences), so $(k)\times(f_{(A,B,k)}=\lp K)$ is contractible.
\end{proof}

\noindent\textbf{Status of the crux.} The Crux below is mathematically
established, through Pres4. What is missing is a proof that avoids that
machinery. The embedding gives \emph{uniqueness} of
deloopings; it does not give their \emph{existence} in the form of a level-$3$
presentation. Precisely, the statement which is proved (via Pres4) but for which a
Pres4-free proof is open is
\begin{statement}[Crux]
For every $t:\Tri$, $\Del\,t\to\Unloop f_t$; equivalently (by
Proposition~\ref{prop:crux} and Theorem~\ref{thm:E}) every $Y:\SCf$ is
$\pteq\Fib K$ for a level-$3$ triple $K$ (a genuine pointed equivalence, not
merely: the triple is unique).
\end{statement}
\noindent It is proved with \code{Presentation4.Pres4} (Proposition~\ref{prop:crux}(3)).
The remaining obligation for this route is a proof of the Crux independent of
Pres4; we do not see one. By (1) it amounts to
showing that $f_t$ agrees with all its translates $D_p$ (the cross term of $k$
vanishes) as soon as $\Fib k$ is a loop space. Classically this holds because
the cross term is the obstruction to $k$ being an $H$-map, and an $H$-structure
on $\Fib k$ over the addition of $K(A,2)$ forces it, the comparison
$\Fib k\wedge\Fib k\to K(A,2)\wedge K(A,2)$ being highly connected; none of this
(smash products, their connectivity, $H$-structures on fibres) is available in
the development; one also needs the classical fact that the primitive
classes in $H^4(K(A,2);B)$ are exactly the suspensions of
$H^5(K(A,3);B)=\mathrm{Hom}(A/2,B)$ (``$H$-map $\Rightarrow$ deloops''). An
equivalent, more internal target: the cross term is a Whitehead square, which
vanishes in an $H$-space. The only existing route to a level-$3$ $K$ is Pres4's: the
fibres of $Y\to\tr{Y}_3$ carry a homogeneous structure (Lemma 4 over the simply
connected $Y$), so $z\mapsto(E_z,\mu_z)$ is a map $\tr Y_3\to K(B,5)$ of
torsors (\code{Classify} at $n=2$). The embedding does not shortcut this: it
compares deloopings of a \emph{given} loop space, while the existence of a
delooping of $\Fib k$ of the form $\Fib K$ is a statement about the cocycle~$k$.

\noindent\textbf{A narrower residual at $t_T$ (suggested by the review; not
checked in detail).} At $A=\Z$: by (H4) (\code{H4Proof.h4}) every
$k:K(\Z,2)\to_*K(B,4)$ is merely $\kappa;K(g,4)$ for a unique $g:\Z\to B$, and
$\Unloop f_t$ is a proposition, so a mere path suffices. If $2g=0$ then
$g=h\circ\rho$ and $k=\bar k;K(h,4)$, so $f_t=\lp(\hat F;K(h,5))$ deloops, by
naturality of $K(h,-)$ and $e_4$ only. Hence at $(\Z,B,k)$ the Crux reduces
to: \emph{if $\Fib(\kappa;K(g,4))$ is a loop space, then $2g=0$} (the cross term
of the cup square is $2g\,\iota\otimes\iota$). For $t_T$ one needs in
addition $\pi_3T\cong\Z$, i.e.\ $\tr T_3\simeq K(\Z,3)$ (Pres4 gives only an
arbitrary $A=\Om^3\tr T_3$). Neither input is the content of Pres4.

\noindent\textbf{Consequences for the assembly.} With (b), M holds for all of
$\SCf$, and \code{Pi4Sphere} is unchanged (it uses M at $T$ and $BR$). The Crux is
a \emph{sufficient} route to a Pres4-free assembly, not a demonstrated
necessity. It is needed only at the single triple $t_T$ of $\Om T$: if
$f_{t_T}$ deloops, to $K$ say, then $T\pteq\Fib K$ by Theorem~\ref{thm:E} and
Proposition~\ref{prop:crux}(2), a genuine pointed equivalence, and (a) gives M at $T$. A second route, which
does not construct any delooping of a cocycle, is to prove M directly for
deloopable targets, in the looped form: $r^*$ bijective into $\real\,t$ for every
deloopable $t$. This is the route on which the embedding could genuinely
simplify the proof, and it is the next thing to investigate. Our analysis
does not show that replacing Pres4 needs another full presentation theorem.

\section{What is formalised}

Module \code{DeloopEmbedding} (spec \code{embedding.opens}; checked together
with the whole development, with no deviation in any \code{.opens} spec):
\begin{itemize}
\item $\SCf$, $\Om$ (\code{SC4}, \code{Acs}, \code{Omega});
\item Lemma~\ref{lem:loopmaps}: \code{SourceS3.LoopEquiv} at $c=2$, used as \code{LoopEq.LE};
\item Lemma~\ref{lem:loopeq}: \code{LoopEq.om}, \code{dl}, \code{dl-eq},
  \code{el}, \code{delta}, \code{om-delta} (by definition), \code{fst-eq},
  \code{loop-delta};
\item Theorem~\ref{thm:E}: \code{E.contr1}, \code{contr2}, \code{contr3};
  \code{Emb.tr-ap}, \code{Emb.emb} ($\ap_\Om$ is an equivalence),
  \code{fib-prop};
\item Corollary~\ref{cor:C}: \code{Deloopable}, \code{deloopable-prop},
  \code{tot-ret}, \code{tot-equiv}, \code{sc4-equiv} $:(t:\Tri)\times\Del\,t\simeq\SCf$;
\item step 1: \code{tbar}, \code{BR-pt}, \code{deloop-R}; \code{Tor-pt},
  \code{deloop-Tor}, \code{BR-Tor} $:BR\pteq\Tor_R$;
\item statement M and (b): \code{Mark[S3,base,surf,ump3].Marked},
  \code{Mark.marked} (every $Y:\SCf$, through Pres4).
\end{itemize}
Not formalised: the looped form of M, Proposition~\ref{prop:crux}, and the Crux
independently of Pres4.
\end{document}
